\subsection{The Timoshenko Trace}\label{sec:trace-timo}
\begin{Prompt}
particular solution:
\begin{equation}
\begin{aligned}
\hat{\bm u}_{\rm P}\{\bm a,\bm{b}_{\Section}\}
&=\Big(-\tfrac12 \reflenx^2\,\mathbf{1}-\nu\,\operatorname{dev}\bm{r}\otimes\bm{r}+\reflenx\,\FrameBasis\otimes\bm{r}\Big)\bm{a}_{\Section}
+(\reflenx\,\FrameBasis-\nu\,\bm{r})\,a_{\varepsilon}
+\big(\reflenx\,\FrameBasis\times\bm{r}+\WarpTwistIesan\,\FrameBasis\big)a_\vartheta\\
&\quad-\tfrac16\,\reflenx^3\,\bm{b}_{\Section}
-\reflenx\,\nu\Big(\operatorname{dev}\bm{r}\otimes\bm{r}-\bm{r}\otimes\CentroidLocation\Big)\bm{b}_{\Section}
+\tfrac12 \reflenx^2\,\FrameBasis\otimes(\bm{r}-\CentroidLocation)\,\bm{b}_{\Section}
+\FrameBasis\otimes\bm{\WarpShearIesan}(\bm{r})\,\bm{b}_{\Section}\\
&\quad+\big(\reflenx\,\FrameBasis\times\bm{r}+\WarpTwistIesan\,\FrameBasis\big)c_\vartheta, 
\end{aligned}
\end{equation}


The infinitesimal strain is
\[
\StrainTensor
= \varepsilon\,\FrameBasis\otimes\FrameBasis
\;+\;\FrameBasis\stackrel{\rm s}{\otimes} \bm{\varepsilon}_{\gamma}
\;-\;\nu\,\varepsilon\,\oneS,
\]
where
\begin{subequations}
\begin{align}
\PointAxialStrain
&= a_{\varepsilon} + \bm{r}\cdot\bm{a}_{\Section}
+ x\,(\bm{r}-\CentroidLocation)\cdot\bm{b}_{\Section}, \\
\PointShearStrain &= (a_{\vartheta} + c_{\vartheta})(\FrameBasis\times \bm{r} + \nabla\WarpTwistIesan) - \nu \bigl(\operatorname{dev}\bm{r}\otimes\bm{r}
- \bm{r}\otimes\CentroidLocation\bigr)\bm{b}_{\Section}
+ (\nabla \bm{\WarpShearIesan})^T \bm{b}_{\Section}
\end{align}
\label{eq:iesan-point-strain}
\end{subequations}

Dont assume any special coordinates. $r_0$ is arbitrary, nothing principal or centroidal.
\end{Prompt}

We work with two orthonormal transverse directions 
$\{\bm{t}_1,\bm{t}_2\}$ in the section plane, with
\[
\FrameBasis\cdot\bm{t}_\alpha = 0,
\qquad
\bm{t}_1\cdot\bm{t}_2 = 0,
\]
and define bending axes
\[
\bm{n}_\alpha := \FrameBasis\times\bm{t}_\alpha,
\qquad \alpha=1,2.
\]
Let $\bm{r}_0\in\Omega$ be a designated reference point (in 2D Timoshenko: 
the neutral axis point $y=0$). The center of twist is denoted $\bm{r}_\vartheta$.

\paragraph{Definition.}

The trace translation is defined pointwise at $\bm{r}_0$:
\begin{equation}
\bm{u}(x) := \hat{\bm{u}}(x,\bm{r}_0).
\label{eq:def-tim-trace-u}
\end{equation}

The bending rotation components about $\bm{n}_\alpha$ are defined by 
Timoshenko-type directional derivatives of the axial displacement:
\begin{equation}
\theta_\alpha(x)
:= \bm{n}_\alpha\cdot\bm{\theta}(x)
:= -\,\bm{t}_\alpha\cdot\nabla_{\bm{r}}
\big(\hat{\bm{u}}(x,\bm{r})\cdot\FrameBasis\big)\big|_{\bm{r}=\bm{r}_0},
\qquad \alpha=1,2.
\label{eq:def-tim-theta-bend}
\end{equation}
The twist component is defined as rotation about the center of twist:
\begin{equation}
\theta_\vartheta(x)
:= \FrameBasis\cdot\bm{\theta}(x)
:= \frac{1}{2}\,\big(\nabla_{\bm{x}}\times\hat{\bm{u}}(x,\bm{r})\big)
\cdot\FrameBasis\Big|_{\bm{r}=\bm{r}_\vartheta}.
\label{eq:def-tim-theta-twist}
\end{equation}
As usual, the full rotation vector is
\begin{equation}
\bm{\theta}(x)
=
\theta_\vartheta(x)\,\FrameBasis
+ \theta_1(x)\,\bm{n}_1 
+ \theta_2(x)\,\bm{n}_2.
\label{eq:def-tim-theta-vector}
\end{equation}

We collect the constant-in-$x$ strain coefficients as
\begin{equation}
\bm{e}_0
=
\begin{pmatrix}
\epsilon \\[2pt]
\bm{\gamma}_{\perp,0} \\[2pt]
\kappa_\vartheta \\[2pt]
\bm{\kappa}_{\perp,0}
\end{pmatrix}
\in\mathbb{R}^6,
\qquad
\bm{e}_0(x;\bm{p})
=
\bm{e}_0 + x\,\bm{e}_1.
\end{equation}


We write the Ieşan parameter vector as
\begin{equation}
\bm{p}
=
\begin{pmatrix}
a_\varepsilon \\[2pt]
\bm{a}_\Omega \\[2pt]
a_\vartheta \\[2pt]
\bm{b}_{\Section}
\end{pmatrix}
\in\mathbb{R}^{1+2+1+2}.
\end{equation}
The SV displacement is decomposed into four modes:
\[
\hat{\bm{u}}
=
\hat{\bm{u}}^{(\varepsilon)} 
+ \hat{\bm{u}}^{(a)} 
+ \hat{\bm{u}}^{(\vartheta)} 
+ \hat{\bm{u}}^{(b)},
\]
corresponding to extension, bending, torsion, and flexure. We now summarize 
the effect of each mode on $(\bm{u},\bm{\theta})$ and hence on the constant 
strain coefficients; detailed SV expressions are identical to those used in 
the section-average derivation and are not repeated here.

\paragraph{Extension.}

For the extension mode
$\hat{\bm{u}}^{(\varepsilon)} = (x\,\FrameBasis - \nu\,\bm{r})\,a_\varepsilon$,
we obtain
\begin{equation}
\bm{u}^{(\varepsilon)}(x) = x\,\FrameBasis\,a_\varepsilon - \nu\,\bm{r}_0\,a_\varepsilon,
\qquad
\bm{\theta}^{(\varepsilon)}(x) = \mathbf{0}.
\end{equation}
Thus
\begin{equation}
\epsilon^{(\varepsilon)} = a_\varepsilon,
\quad
\bm{\gamma}_{\perp,0}^{(\varepsilon)} = \mathbf{0},
\quad
\kappa_\vartheta^{(\varepsilon)} = 0,
\quad
\bm{\kappa}_{\perp,0}^{(\varepsilon)} = \mathbf{0}.
\end{equation}

\paragraph{Bending.}

For the bending mode $\hat{\bm{u}}^{(a)}$, the axial displacement is
\[
\hat{\bm{u}}^{(a)}(x,\bm{r})\cdot\FrameBasis
=
-\tfrac{1}{2}x^2\,(\bm{a}_\Omega\cdot\bm{r}_\perp)
- \nu\,\operatorname{dev}(\bm{r}_\perp\otimes\bm{r}_\perp)\bm{a}_\Omega
+ x\,(\bm{r}_\perp\cdot\bm{a}_\Omega),
\]
where $\bm{r}_\perp$ is the in-plane projection of $\bm{r}$. Applying the 
Timoshenko bending rotation definition \eqref{eq:def-tim-theta-bend} at 
$\bm{r}=\bm{r}_0$ yields
\begin{equation}
\theta_\alpha^{(a)}(x)
=
x\,(\bm{t}_\alpha\cdot\bm{a}_\Omega)
+ \Xi_{\alpha}^{(a)}\,\bm{a}_\Omega,
\qquad
\alpha=1,2,
\end{equation}
where the $2\times 2$ tensor
\begin{equation}
\Xi^{(a)} \in \mathbb{R}^{2\times 2},\quad
(\Xi^{(a)})_{\alpha\beta}
:=
\nu\,\bm{t}_\alpha\cdot 
\big[\operatorname{dev}(\bm{r}_0\otimes\bm{r}_0)\,\bm{e}_\beta\big],
\end{equation}
encodes the Poisson-induced bending rotation at $\bm{r}_0$, and 
$\{\bm{e}_\beta\}$ is the in-plane basis aligned with $\bm{a}_\Omega$ 
components. The twist component from pure bending remains zero:
\begin{equation}
\theta_\vartheta^{(a)}(x) = 0.
\end{equation}

Direct differentiation gives the curvature contributions
\begin{equation}
\bm{\kappa}_{\perp,0}^{(a)}
=
\bm{\theta}'^{(a)}(0)
=
\sum_{\alpha=1}^2 (\bm{t}_\alpha\cdot\bm{a}_\Omega)\,\bm{n}_\alpha
=
-\,\FrameBasis\times\bm{a}_\Omega,
\qquad
\kappa_\vartheta^{(a)} = 0.
\end{equation}
The bending contribution to $\bm{u}(x)$ differentiates to
\[
\bm{u}'^{(a)}(x)
=
- x\,\bm{a}_\Omega,
\]
so that the shear strain from bending is
\begin{equation}
\bm{\gamma}_{\perp,0}^{(a)}
=
\FrameBasis\times(\Xi^{(a)}\bm{a}_\Omega).
\end{equation}

\paragraph{Torsion.}

For the torsion mode 
$\hat{\bm{u}}^{(\vartheta)} = (x\,\FrameBasis\times\bm{r} + \omega(\bm{r})\,\FrameBasis)\,a_\vartheta$,
the trace evaluates rotation about the center of twist $\bm{r}_\vartheta$:
\begin{equation}
\theta_\vartheta^{(\vartheta)}(x)
=
\big(x\,\FrameBasis + \bm{c}_\vartheta\big)\cdot\FrameBasis\,a_\vartheta
=
x\,a_\vartheta + c_\vartheta\,a_\vartheta,
\end{equation}
where $\bm{c}_\vartheta = \bm{r}_\vartheta$ and $c_\vartheta$ is the 
flexure-induced twist offset (zero in pure torsion). The bending components 
remain zero:
\begin{equation}
\theta_\alpha^{(\vartheta)}(x) = 0,\quad \alpha=1,2.
\end{equation}
Thus
\begin{equation}
\kappa_\vartheta^{(\vartheta)} = a_\vartheta,
\qquad
\bm{\kappa}_{\perp,0}^{(\vartheta)} = \mathbf{0},
\qquad
\bm{\gamma}_{\perp,0}^{(\vartheta)} = \FrameBasis\times(\bm{c}_\vartheta\,a_\vartheta).
\end{equation}

\paragraph{Flexure.}

For the flexure mode $\hat{\bm{u}}^{(b)}$, the Timoshenko rotations are
\begin{equation}
\theta_\alpha^{(b)}(x)
=
-\,\bm{t}_\alpha\cdot\nabla_{\bm{r}}
\big(\hat{\bm{u}}^{(b)}(x,\bm{r})\cdot\FrameBasis\big)\big|_{\bm{r}=\bm{r}_0},
\qquad \alpha=1,2.
\end{equation}
Splitting $\hat{\bm{u}}^{(b)}$ into the classical polynomial SV part and the 
shear warping contributions yields
\begin{equation}
\bm{\theta}^{(b)}(x)
=
-\tfrac{1}{2}x^2\,(\FrameBasis\times\bm{b}_{\Section})
+ x\,\Xi^{(b)}_\theta\,\bm{b}_{\Section}
+ \Xi^{(\psi)}_\theta\,\bm{b}_{\Section},
\end{equation}
with
\begin{align}
\Xi^{(b)}_\theta &\in\mathbb{R}^{3\times 2},
&(\Xi^{(b)}_\theta)_{\bullet\beta} 
&:= \sum_{\alpha=1}^2 \bm{n}_\alpha\,
\Big(-\,\bm{t}_\alpha\cdot\big[\bm{W}_{(b)}(\bm{r}_0)\,\bm{e}_\beta\big]\Big),
\\[2pt]
\Xi^{(\psi)}_\theta &\in\mathbb{R}^{3\times 2},
&(\Xi^{(\psi)}_\theta)_{\bullet\beta}
&:= \sum_{\alpha=1}^2 \bm{n}_\alpha\,
\Big(-\,\bm{t}_\alpha\cdot[\nabla_{\bm{r}}\bm{\psi}(\bm{r}_0)\,\bm{e}_\beta]\Big),
\end{align}
collecting the contributions from the polynomial flexure field and the shear 
warping, respectively. Differentiation yields
\begin{equation}
\bm{\kappa}_{\perp,0}^{(b)}
=
\bm{\theta}'^{(b)}(0)
=
\big(\FrameBasis\otimes\bm{r}_s + \Xi^{(b)}_\theta\big)\bm{b}_{\Section},
\end{equation}
where $\bm{r}_s$ is the shear center. The twist curvature contribution is
\begin{equation}
\kappa_\vartheta^{(b)}
=
\partial_{\bm{b}}c_\vartheta\cdot\bm{b}_{\Section},
\end{equation}
as in the standard SV flexure–torsion coupling.

The translation from flexure gives
\begin{equation}
\bm{u}'^{(b)}(0)
=
-\,\tfrac{1}{2}\,\bm{b}_{\Section}
+ \Gamma^{(b)}_0\,\bm{b}_{\Section},
\end{equation}
for a suitable $2\times 2$ tensor $\Gamma^{(b)}_0$ collecting the dependence 
on $\bm{r}_0$ and warping values. Thus the shear strain contribution is
\begin{equation}
\bm{\gamma}_{\perp,0}^{(b)}
=
\big(-\tfrac{1}{2}\mathbf{1}_\Omega + \Gamma^{(b)}_0\big)\bm{b}_{\Section}
+ \FrameBasis\times\big(\Xi^{(b)}_\theta\bm{b}_{\Section}\big).
\end{equation}

\paragraph{Assembly of the trace matrices.}

Collecting the constant strain coefficients from all four modes, we obtain
\begin{equation}
\bm{e}_0(\bm{p})
=
\TraceDeDpRoot\,\bm{p},
\qquad
\bm{e}_1(\bm{p})
=
\TraceDeDpOne\,\bm{p},
\end{equation}
with $\bm{e}_0 = (\epsilon,\bm{\gamma}_{\perp,0},\kappa_\vartheta,
\bm{\kappa}_{\perp,0})^\top$ and $\bm{p}=(a_\varepsilon,\bm{a}_\Omega,
a_\vartheta,\bm{b}_{\Section})^\top$. The Timoshenko–torsion hybrid trace matrix 
has the block form
\begin{equation}
\TraceDeDpRoot
=
\begin{pmatrix}
1 & \mathbf{0}^\top & 0 & 0 \\[4pt]
\mathbf{0} 
& \big(\FrameBasis\times\Xi^{(a)}\big)_\perp 
& \big(\FrameBasis\times\bm{c}_\vartheta\big)_\perp 
& \Big(-\tfrac{1}{2}\mathbf{1}_\Omega 
      + \Gamma^{(b)}_0 
      + \FrameBasis\times\Xi^{(b)}_\theta\Big)_\perp \\[4pt]
0 & 0 & 1 & \partial_{\bm{b}}c_\vartheta \\[4pt]
\mathbf{0} 
& -(\FrameBasis\times)_\perp 
& \mathbf{0} 
& \big(\FrameBasis\otimes\bm{r}_s + \Xi^{(b)}_\theta\big)_\perp
\end{pmatrix},
\label{eq:T0-tim-hybrid}
\end{equation}
and the universal first-order matrix $\TraceDeDpOne$ (capturing the $x$-linear 
coefficients) has the same structure as in the general SV analysis, namely
\begin{equation}
\TraceDeDpOne
=
\begin{pmatrix}
0 & \mathbf{0}^\top & 0 & 0 \\[4pt]
\mathbf{0} & \mathbf{0}^\top & 0 & \Gamma^{(b)}_1 \\[4pt]
0 & \mathbf{0}^\top & 0 & 0 \\[4pt]
\mathbf{0} & -\FrameBasis^{\times} & 0 & -\FrameBasis\times\mathbf{1}_\Omega
\end{pmatrix},
\label{eq:T1-tim-hybrid}
\end{equation}
where $\Gamma^{(b)}_1$ collects the $x$-linear shear gradient coefficients 
from the flexure mode (unchanged across SV-equivalent traces). 

In 2D plane bending with a single transverse direction $\bm{t}_1$, a single 
bending axis $\bm{n}_1$, and torsion suppressed, the above reduces exactly to 
Timoshenko’s original pointwise rotation definition 
$\theta(x) = -\partial\hat{u}_x/\partial y|_{y=y_0}$ and reproduces the 
classical $2/3$ shear factor.
